% Auto-generated by generate_tex_fragments.py -- do not hand-edit.
\begin{longtable}{@{}L{2.0cm}L{2.0cm}L{2.1cm}L{6.2cm}L{1.9cm}@{}}
\caption{P4: A stop rule prevents routine joint education from worsening risk when fear or coercion appears. Record, design, population/setting, finding as charted (verbatim from the library record), and direction relative to the proposition. n=23 linked records (kg/graph.json edges).}\label{tab:p4}\\
\toprule
\textbf{Record} & \textbf{Design} & \textbf{Population\slash{}setting} & \textbf{Finding (as charted)} & \textbf{Direction}\\
\midrule
\endfirsthead
\multicolumn{5}{l}{\textit{P4 continued}}\\
\toprule
\textbf{Record} & \textbf{Design} & \textbf{Population\slash{}setting} & \textbf{Finding (as charted)} & \textbf{Direction}\\
\midrule
\endhead
\bottomrule
\endfoot
Couture-Carron A 2020 \cite{srp31896308} & \seqsplit{Exploratory} \seqsplit{qualitative} study & South Asian Muslim \seqsplit{communities} (diaspora) & Cultural norms \seqsplit{constructing} dating as shameful were found to intensify fear of parental\slash{}community reaction and strengthen attachment to abusive partners, with direct \seqsplit{implications} for disclosure, help-seeking, and the ability to end abusive \seqsplit{relationships}. & Supports\\
\\[3pt]
Han X 2024 \cite{srp38970860} & Cross-sectional, random-effects Poisson regression (n=201) & Married Salar (Turkic-Muslim minority) women, rural China & Among 201 married Salar Muslim women, 65.6\% reported past-year child neglect, 30.8\% reported past-year IPV, and 37.6\% had \seqsplit{experienced} child marriage; child marriage and IPV \seqsplit{victimization} were associated with downstream child-neglect risk. & Supports\\
\\[3pt]
Hegarty K 2025 \cite{srp41330529} & \seqsplit{Qualitative} directed content analysis of open-ended survey responses (n=682 women with lived experience of coercive control) & Australian women who had \seqsplit{experienced} coercive control & Findings align with the WHO LIVES framework (Listen-Inquire-Validate-Enhance safety-Support) and the CARE framework (Choice and control-Action and advocacy-\seqsplit{Recognition} and \seqsplit{understanding}-Emotional connection); women wanted specific validating, non-judgmental scripted language from \seqsplit{practitioners} and messaging that \seqsplit{counteracts} coercive-\seqsplit{controlling} tactics rather than generic \seqsplit{acknowledgement}. (Verifier correction recorded in the library entry.) & Supports\\
\\[3pt]
Islam MK 2021 \cite{srp34527969} & Mixed methods: cross-sectional survey (n=96) + \seqsplit{qualitative} interviews (n=18) + provider interviews (n=9) & Rohingya adolescent girls (age 10-19) with experience of child marriage, Cox's Bazar refugee camp, Bangladesh & Mean age at first marriage was 15.7 years; over 80\% had given birth or were pregnant within two years of the survey; key drivers of child marriage included perceived physical\slash{}mental maturity norms, social norms, insecurity, and family honour. & Supports\\
\\[3pt]
James LE 2021 \cite{srp34770188} & Co-designed, two-phase community-based \seqsplit{intervention} (workshops); RCT-labelled but abstract available describes \seqsplit{intervention} \seqsplit{development}\slash{}testing & Rohingya (Muslim minority, displaced) in Malaysia and Syrian refugees (Muslim-majority) in Lebanon & A social-norms-based, mental-health-integrated community \seqsplit{intervention} was co-designed with Rohingya and Syrian refugee \seqsplit{communities} \seqsplit{specifically} to reduce IPV and encourage help-seeking, using culturally embedded workshop delivery rather than clinical\slash{}individual counseling alone. & Supports\\
\\[3pt]
Kalichman SC 2020 \cite{srp32869478} & Commentary \slash{} research-agenda-setting article & Adolescent girls and young women (15-24) in LMICs receiving HIV index-testing\slash{}partner-\seqsplit{notification} services & Highlights that partner-\seqsplit{notification} and index-testing approaches proven safe\slash{}effective for adult men and cohabiting\slash{}married adults have not been adequately studied for younger women with less \seqsplit{relationship} power, who face higher risk of IPV, stigma, and economic coercion when disclosure is required or \seqsplit{facilitated} by a health service. (Verifier correction recorded in the library entry.) & Supports\\
\\[3pt]
Karakurt 2016 \cite{ref45} & systematic review and meta-analysis (6 studies of 1,733 screened citations) & couples in treatment for physical\slash{}\seqsplit{situational} IPV, mostly US clinical\slash{}research samples & After screening 1,733 citations, only 6 studies met inclusion criteria for meta-analysis of couples therapy as a treatment for IPV; \seqsplit{preliminary} pooled data suggest couples therapy can be a viable treatment 'in select situations' but the evidence base is small and \seqsplit{professionals} remain cautious about risk of violent \seqsplit{retaliation} between partners. & Supports\\
\\[3pt]
Kuswanto H 2024 \cite{ref3} & Cross-sectional, binary logistic regression (nationally \seqsplit{representative}, n=9,333) & Females aged 15-20, Indonesia (2017 Indonesia \seqsplit{Demographic} and Health Survey) & National prevalence of female child marriage in Indonesia was 12.53\%; health insurance status and sex of household head were not \seqsplit{significant} predictors, while other \seqsplit{socioeconomic}\slash{}education factors were associated with child marriage risk. & Supports\\
\\[3pt]
McCauley HL 2017 \cite{ref17} & Secondary \seqsplit{psychometric} analysis of pooled \seqsplit{multicenter}\slash{}RCT survey data & 4,674 women aged 16-29 in 24 \seqsplit{Pennsylvania} and 5 California family planning clinics, USA & Confirmed a 2-domain structure (pregnancy coercion, condom \seqsplit{manipulation}) and produced a validated 5-item short-form \seqsplit{Reproductive} Coercion Scale (RCS); recent \seqsplit{reproductive} coercion prevalence was 6.3-6.7\% depending on the scale version used. & Supports\\
\\[3pt]
McCollum 2008 \cite{ref46} & narrative review of outcome research and clinical practice & IPV-affected couples in conjoint treatment settings, US clinical literature & Concludes conjoint IPV treatment can be effective and safe only when embedded in a broader community response, with careful screening of couples for inclusion, \seqsplit{modification} of standard conjoint methods to promote safety, and ongoing safety assessment with \seqsplit{contingency} plans for escalation; a 'one-size-fits-all' approach is explicitly rejected given IPV typologies. & Supports\\
\\[3pt]
Mieth K 2025 \cite{srp40091116} & \seqsplit{Qualitative} study (49 in-depth interviews, 16 focus group \seqsplit{discussions}) & Forcibly Displaced Myanmar Nationals (Rohingya) aged 15-24, Cox's Bazar, Bangladesh (post-October 2016 arrivals) & \seqsplit{Participants} reported perceived increases in child marriage rates post-\seqsplit{displacement}, attributed to protection fears, \seqsplit{socioeconomic} need, lost education\slash{}employment \seqsplit{opportunity}, and perceived loosening of legal-age \seqsplit{enforcement} in the camp context. & Supports\\
\\[3pt]
Nelson HD 2012 \cite{srp22565034} & systematic review (evidence base for USPSTF 2013 update) & Women in health care settings, US-relevant evidence base & Of 15 fair\slash{}good-quality studies evaluating 13 screening \seqsplit{instruments}, 6 were highly accurate; 14 studies found minimal adverse effects from screening, though some women reported discomfort, loss of privacy, emotional distress, and fear of further abuse; counseling \seqsplit{interventions} reduced IPV\slash{}improved outcomes in some \seqsplit{subpopulations} (pregnant women, new mothers, family-planning clients) but a large trial found no \seqsplit{significant} screening-vs-usual-care difference overall. & Supports\\
\\[3pt]
Onukwugha FI 2025 \cite{ref42} & Quasi-\seqsplit{experimental} matched-pair cluster-controlled trial (n=1617 clients, 20 \seqsplit{intervention} vs 20 comparison clinics) & Family planning and antenatal care clients, Nigeria & Combined WHO LIVES + ARCHES first-line response reduced relative risk of IPV exposure (beta=0.54, 95\% CI 0.30-0.98) and \seqsplit{reproductive} coercion (beta=0.51, 95\% CI 0.28-0.94) over 9 months among women not already \seqsplit{experiencing} violence, but showed no effect on modern \seqsplit{contraceptive} use; authors note LIVES and ARCHES \seqsplit{individually} have shown 'mixed effects' in prior trials. (Verifier correction recorded in the library entry.) & Supports\\
\\[3pt]
Rumble L 2018 \cite{ref4} & Cross-sectional secondary analysis, \seqsplit{multivariate} regression (nationally \seqsplit{representative}, n=6,578) & Women aged 20-24, Indonesia (2012 Indonesian \seqsplit{Demographic} and Health Survey \slash{} Adolescent \seqsplit{Reproductive} Health Survey) & In a nationally \seqsplit{representative} Indonesian sample, \seqsplit{approximately} 17\% of women were married before age 18 and 6\% before age 16; the paper notes research on child marriage in Southeast Asia \seqsplit{specifically} was scarce prior to this nationally \seqsplit{representative} \seqsplit{multivariate} analysis. & Supports\\
\\[3pt]
Williamson HC 2015 \cite{ref10} & RCT (N=130 engaged\slash{}newlywed couples randomized to 1 of 3 \seqsplit{relationship}-education programs, 3-year follow-up) & US engaged and newlywed couples & Couples with acute baseline concerns, including higher physical aggression and alcohol use, benefited LESS from \seqsplit{relationship}-education \seqsplit{intervention} than couples without these concerns; couples with high baseline \seqsplit{satisfaction}\slash{}commitment actually declined faster in \seqsplit{satisfaction} under an intensive skill-based \seqsplit{intervention} versus a low-intensity one. & Supports\\
\\[3pt]
Wilson 2023 \cite{ref99} & instrument-\seqsplit{development}\slash{}validation study & adult \seqsplit{participants} reporting on intimate \seqsplit{relationships} (validation sample) & Reports on validity\slash{}\seqsplit{reliability} of the Demand, Threat, \seqsplit{Surveillance}, and Response-to-Demands subscales of the Coercion in Intimate Partner \seqsplit{Relationships} (CIPR) scale, developed because existing coercive-control measures had documented content-validity gaps (not capturing the full range of demand\slash{}threat\slash{}\seqsplit{surveillance} tactics). & Supports\\
\\[3pt]
Zust BL 2021 \cite{ref118} & Repeated cross-sectional survey (\seqsplit{convenience} sample of Upper-Midwest US \seqsplit{congregations}, 2005 and 2015) & Christian clergy and \seqsplit{congregation} members, US & Faith \seqsplit{commitments} led some domestic-violence victims to endure abuse rather than compromise their faith by leaving; over 10 years, \seqsplit{congregational} support for victims changed slowly; clergy wanted more training in counseling and speaking about domestic violence, and felt \seqsplit{comfortable} making referrals but not \seqsplit{necessarily} providing direct DV counseling themselves. (Verifier correction recorded in the library entry.) & Supports\\
\\[3pt]
Markman HJ 2013 \cite{ref9} & Randomized controlled trial (couples via religious \seqsplit{organizations} randomized to PREP vs. naturally occurring premarital services; n=44\slash{}66\slash{}83) & US couples marrying through religious \seqsplit{organizations}, 8-year follow-up on divorce & No overall difference in divorce rates between PREP (clergy-delivered or \seqsplit{professional}-delivered) and usual premarital services; couples with higher premarital negative-\seqsplit{communication} observed pre-\seqsplit{intervention} were MORE likely to divorce if given PREP, while low-negativity couples benefited; a premarital history of physical aggression showed a similar (marginally \seqsplit{significant}) harmful-moderation pattern. & Complicates\\
\\[3pt]
Miller Elizabeth 2016 \cite{ref110} & Cluster randomized controlled trial (ARCHES) & 4,009 females aged 16-29 across 25 family planning clinics (17 clusters), USA & The universal education\slash{}counseling \seqsplit{intervention} (ARCHES) did NOT \seqsplit{significantly} reduce \seqsplit{reproductive} coercion (ARR 1.50, 95\% CI 0.95-2.35) or partner violence (ARR 1.07, 95\% CI 0.84-1.38) at 1-year follow-up, though it improved secondary outcomes such as \seqsplit{recognition} and self-efficacy. & Complicates\\
\\[3pt]
Robertson 2011 \cite{ref49} & cross-sectional \seqsplit{comparative} study, n=172 (85 men, 87 women), self-report (CTS2 + modified PMWI) & community\slash{}\seqsplit{convenience} samples reporting on own and partner's IPV behavior & Coercive control was associated with IPV similarly for men and women across three samples, and was \seqsplit{predominantly} reciprocal --- both partners commonly reported both receiving and \seqsplit{perpetrating} \seqsplit{controlling} behaviors, regardless of physical-abuse status. & Complicates\\
\\[3pt]
Schneider 2014 \cite{ref47} & \seqsplit{theoretical}\slash{}clinical article (non-empirical) & couples presenting for therapy with \seqsplit{situational} couple violence (SCV) & Argues conjoint treatment can be \seqsplit{appropriate} for \seqsplit{situational} couple violence \seqsplit{specifically} (as distinct from coercive \seqsplit{controlling} violence), proposing an attachment-based safety plan and time-out strategy, while explicitly flagging cautions about \seqsplit{appropriateness} of couples counseling when violence is ongoing. & Complicates\\
\\[3pt]
Portnoy 2020 \cite{srp32791967} & \seqsplit{qualitative} study --- patient interviews (n=10) and provider focus groups (n=29) & US Veterans Health \seqsplit{Administration} patients and providers & Identifies convergent and divergent patient\slash{}provider beliefs about barriers and \seqsplit{facilitators} to screening \seqsplit{specifically} for IPV \seqsplit{perpetration} (not just \seqsplit{victimization}), an area with a much thinner evidence base than \seqsplit{victimization} screening. & Relevant only\\
\\[3pt]
Stern E 2018 \cite{ref114} & \seqsplit{Implementation}\slash{}adaptation review (\seqsplit{qualitative} evaluation and programme \seqsplit{documentation}) & \seqsplit{Indashyikirwa} programme, \seqsplit{implemented} by CARE Rwanda, RWAMREC and RWN, rural Rwanda & \seqsplit{Indashyikirwa} combined a 5-month couples curriculum, sustained community activism by trained couples, women's safe spaces for dedicated support and referral of IPV survivors, and training of opinion leaders to support an enabling \seqsplit{environment}; adaptation was guided by the SASA! fidelity brief. & Relevant only\\
\\[3pt]
\end{longtable}
